% =============================================================================
% Privacy edit — accessibility analytics.
% The analytics event now records a SINGLE non-identifying boolean,
% "accessibilityUsed" (whether ANY accessibility preference was active), and
% never which specific mode. Update the paper text that still lists the modes.
% Implemented in LLM_Call.log_answer_completed (the answer_completed event).
% =============================================================================


% -----------------------------------------------------------------------------
% PAPER A ("Renaixement Museum Interactive Audioguide" version) —
% "Metadata-Only Analytics" subsection.
% In the enumerate of collected fields, REPLACE the last item
%   "\item response time, selected language, room, artwork, and enabled
%    accessibility modes."
% with:
% -----------------------------------------------------------------------------

    \item response time, selected language, room, and artwork;
    \item whether any accessibility option was active, recorded as a single
          yes/no flag and never which specific option.

% And, in the paragraph after the list, the sentence
%   "The text of visitor questions and generated answers is not stored in this
%    analytics file."
% can be extended to:

% The text of visitor questions and generated answers is not stored in this
% analytics file, and the specific accessibility modes a visitor selects are not
% stored either. Only a single non-identifying flag indicates whether any
% accessibility option was active, so the museum can observe overall uptake of
% these features without recording data that could reveal an individual
% visitor's disability or specific needs.


% -----------------------------------------------------------------------------
% PAPER B ("Towards Adaptive and Situated AI..." version) —
% "Administrative Interface Design" subsection.
% The sentence
%   "This includes navigation paths, frequently accessed artworks, common
%    queries, and selected language preferences."
% can be extended (optional, only if you want to make the privacy stance explicit):
% -----------------------------------------------------------------------------

% This includes navigation paths, frequently accessed artworks, common queries,
% and selected language preferences. To protect sensitive information, the
% specific accessibility options a visitor selects are not recorded; the system
% stores only a single non-identifying flag indicating whether any accessibility
% option was used, allowing the museum to gauge uptake of these features without
% revealing an individual visitor's needs.
